\paragraph{Complete recognition.}
The controlled follow-up uses nineteen fresh source inputs, eight arms
(old and current defaults and batch portfolios, each with an A/A copy),
one warm-up and five shuffled measured rounds. All 760 measured calls and
152 warm-ups complete; the run takes 84.716 seconds. Timing includes the
entire recognition call and mandatory source replay. Ratios are medians
of paired old/new times, not quotients of the marginal time medians.

\begin{center}\small
\begin{tabular}{@{}lrrrrr@{}}
\toprule
Input & default ms & old port. ms & new port. ms & old/new & default/new \\
\midrule
survivor-00 & 17.192 & 15.347 & 15.811 & 0.993 & 1.115 \\
survivor-01 & 17.096 & 46.533 & 41.473 & 0.967 & 0.487 \\
survivor-02 & 15.871 & 10.243 & 10.323 & 0.992 & 1.527 \\
survivor-03 & 11.346 & 43.427 & 43.755 & 0.988 & 0.259 \\
mirror-03 & 24.675 & 6.440 & 6.455 & 1.009 & 3.830 \\
survivor-04 & 10.105 & 31.943 & 31.816 & 1.007 & 0.319 \\
survivor-05 & 19.303 & 16.361 & 16.452 & 0.992 & 1.173 \\
survivor-06 & 7.636 & 5.634 & 5.761 & 0.983 & 1.326 \\
survivor-07 & 19.798 & 21.842 & 21.486 & 1.003 & 0.921 \\
survivor-08 & 7.751 & 7.432 & 7.554 & 0.983 & 1.026 \\
mirror-08 & 20.526 & 16.625 & 16.524 & 1.009 & 1.243 \\
survivor-09 & 7.602 & 5.801 & 5.852 & 0.995 & 1.304 \\
survivor-10 & 6.640 & 5.444 & 5.532 & 0.992 & 1.185 \\
survivor-11 & 7.126 & 5.096 & 5.174 & 0.982 & 1.374 \\
gordian & 1518.937 & 1509.908 & 1623.636 & 0.941 & 0.936 \\
trefoil & 0.052 & 0.053 & 0.053 & 1.004 & 0.995 \\
figure\_eight & 0.060 & 0.060 & 0.059 & 1.003 & 1.013 \\
conway & 0.990 & 1.006 & 0.988 & 1.018 & 1.002 \\
kinoshita\_terasaka & 1.028 & 1.023 & 1.019 & 1.002 & 1.003 \\
\bottomrule
\end{tabular}
\end{center}


The default old/current ratios range from 0.980 to 1.023. Old/new portfolio
ratios range from 0.941 to 1.018, so this sample establishes no broad
whole-input improvement. The four A/A ranges are respectively
0.987--1.016, 0.987--1.031, 0.986--1.016 and 0.967--1.038.
Against the current default, the new optional portfolio still ranges from
0.259 on survivor-03 to 3.830 on mirror-03. Those large gains and regressions
come from the alternative elimination schedule, already present in the
baseline, and should not be attributed to this reachability optimization.

Gordian exposes the price of the larger trial. Recovery costs 39,904 work
units, after which the new search spends its continuation allowance and
falls back. Its trial charge is 89,905 rather than 50,001; total producer
work increases from 1,265,146 to 1,305,050. The fallback proof remains the
same 9,663-byte version-five certificate and its own search still uses
9,025 nodes. Whole-call medians are 1,509.908 ms for the old portfolio and
1,623.636 ms for the new one, with a paired ratio of 0.941 and a new-portfolio
A/A ratio of 1.038. The extra speculative work is real even though noise
affects its timing estimate. Optional status and the existing default are
retained.


\paragraph{Checked group stages on larger circles.}
A separate run bypasses earlier simplification and compares the old and new
direct producers and portfolios, each with an A/A copy. Every successful
call includes source reconstruction, production and mandatory compressed
replay. Seven circle diagrams, one warm-up and five shuffled measured rounds
give 280 measured calls and 56 warm-ups in 45.563 seconds. Of these,
270 measurements and 54 warm-ups complete. All ten incomplete measurements
and both incomplete warm-ups are the old portfolio at 512 crossings, where
fallback hits the unchanged 100,000-node cap. They remain in the evidence;
no old/new completion-time ratio is reported for that portfolio case.

\begin{center}\small
\begin{tabular}{@{}lrrrr@{}}
\toprule
Crossings & direct old/new ms & ratio & portfolio old/new ms & ratio \\
\midrule
8 & 1.409/1.430 & 0.977 & 1.363/1.382 & 0.986 \\
16 & 2.920/2.882 & 1.021 & 2.818/2.824 & 1.004 \\
32 & 6.141/5.797 & 1.060 & 5.915/5.629 & 1.051 \\
64 & 13.458/12.217 & 1.128 & 13.218/12.003 & 1.075 \\
128 & 30.760/24.467 & 1.288 & 377.400/23.901 & 15.998 \\
256 & 74.892/49.422 & 1.515 & 1503.725/50.714 & 30.105 \\
512 & 227.287/101.014 & 2.326 & --/97.486 & -- \\
\bottomrule
\end{tabular}
\end{center}


At 256 crossings the direct checked call improves from 74.892 to 49.422 ms,
with paired ratio 1.515. Producer work drops from 157,823 to 94,320 while
both use 2,544 nodes and the same 12,775-byte certificate. At 512 crossings
the direct ratio is 2.326, work drops from 446,591 to 188,528, and both use
5,104 nodes and the same 25,879-byte certificate. The producer removes
255 or 511 generators in one batch. The reduction is entirely in planning:
these chains require no propagation through old ancestors.

The revised host recovers the direct batch gains previously forfeited by
its trial budget. At 128 crossings its paired ratio against the old host
is 15.998; at 256 it is 30.105. At 256 crossings, recovery consumes 72,677
units and subsequent search 21,643, within the new continuation allowance.
At 512 crossings, recovery consumes 145,125 units and search 43,403;
the new host completes in a median 97.486 ms without fallback. The old
host returns no verdict at its node limit, so this last result demonstrates
a finite-budget coverage difference, not a measured speedup factor.

The stage A/A ranges are 0.979--1.019 for old direct, 0.830--1.023 for
new direct, 0.986--1.031 for completed old portfolios and 0.975--1.030 for
new portfolios. In particular the 128-crossing new-direct control is 0.830,
so its smaller direct timing gain should be treated cautiously. These
controlled stages show where the mechanism helps, while the complete-call
corpus above records the cost and absence of a broad improvement. Neither
set establishes an asymptotic bound for general knot recognition.
